% answers_short.tex —— 1959 年·全国卷 · 简短答案速查
% 由 tools/make_short_answers.py 自 answers.tex 的 \gksolhead/\gkanswer 提取，
% 并与 1959/1959_数学真题_答案要点.md 逐题比对后生成，勿手改（改源 MD 后重跑）。

\begin{gknotice}
\textbf{使用说明}：本节只给最终答案便于快速核对，完整推导见解析版。答案为编写者独立推导，\textbf{非官方标准答案}。
\end{gknotice}

\gkpart{小题}
\begin{itemize}[leftmargin=2.6em]
  \item[\textbf{1.}] $1.5441$
  \item[\textbf{2.}] $-2$
  \item[\textbf{3.}] $-3<x<\dfrac12$
  \item[\textbf{4.}] $-\dfrac{\sqrt6+\sqrt2}{4}$
  \item[\textbf{5.}] $150\pi\ \mathrm{cm}^2$
  \item[\textbf{6.}] 四面体 $ABCD$ 的体积与 $C$、$D$ 在各自平行线上所取的位置无关
\end{itemize}

\gkpart{大题}
\begin{itemize}[leftmargin=2.6em]
  \item[\textbf{7.}] $\left(\frac14,\frac54,\frac94\right)\ \text{\text{或}}\ (1,5,9)$
  \item[\textbf{8.}] $(AB,BC)=(\sqrt7+\sqrt3,\sqrt7-\sqrt3)$ 或 $(\sqrt7-\sqrt3,\sqrt7+\sqrt3)$
  \item[\textbf{9.}] （1）$\sin\angle PBA=\dfrac2{\sqrt5}$，$\cos\angle PBA=\dfrac1{\sqrt5}$，$\tan\angle PBA=2$；（2）$PB=\dfrac a{\sqrt5}$；（3）距离为 $\dfrac{2a}{5}$
  \item[\textbf{10.}] $FG=FE$
\end{itemize}
